\section{Background}


\subsection{Content-Addressable Storage}

Content-addressable storage (CAS) identifies data by its cryptographic hash rather than its location. First formalized by the InterPlanetary File System (\IPFS{})~\cite{benet2014ipfs}, CAS provides:

\begin{definition}[Content Identifier]
A Content Identifier (CID) is a tuple $\langle \text{codec}, \text{multihash} \rangle$ where:
\begin{itemize}
    \item \texttt{codec} specifies the data encoding (e.g., dag-cbor, raw)
    \item \texttt{multihash} is $\langle \text{algo}, \text{len}, \text{digest} \rangle$ identifying the hash function and digest
\end{itemize}
\end{definition}

For example, CID \texttt{bafybeigdyrzt5sfp7udm7hu76uh7y26nf3efuylqabf3oclgtqy55fbzdi} uniquely identifies a specific data object globally.

\subsection{Merkle Trees for Verification}

Merkle trees enable efficient membership proofs:

\begin{theorem}[Merkle Proof Efficiency]
Given a set $S$ of $n$ elements stored in a Merkle tree, verifying membership of element $e \in S$ requires only $O(\log n)$ hash computations and $O(\log n)$ proof size.
\end{theorem}

This property is essential for verifying experience integrity without downloading the entire library.

\subsection{Semantic Memory in AI}

Prior work on AI memory systems includes:

\begin{itemize}
    \item \textbf{Episodic Memory}~\cite{graves2016hybrid}: Neural networks with explicit memory modules
    \item \textbf{Retrieval-Augmented Generation}~\cite{lewis2020retrieval}: Combining retrieval with generation
    \item \textbf{Vector Databases}~\cite{pinecone2021}: Storing embeddings for similarity search
\end{itemize}

These approaches focus on single-agent memory. Our work extends to \emph{distributed} memory shared across agents.

